\documentclass{article}

% 截至 2026-03-15，NeurIPS 2026 官方模板尚未公开；
% 因此这里使用最新可获得的官方模板 neurips_2025.sty。
\usepackage[dandb,preprint]{neurips_2025}

\usepackage[UTF8]{ctex}
\usepackage{hyperref}
\usepackage{url}
\usepackage{booktabs}
\usepackage{amsfonts}
\usepackage{amsmath,amssymb}
\usepackage{nicefrac}
\usepackage{microtype}
\usepackage{xcolor}
\usepackage{graphicx}
\usepackage{enumitem}

\title{模型究竟看见了什么？\\
Ordinary-Bench：通过受控探测与场景重建评估视觉语言模型的空间信念}

\author{匿名作者}

\begin{document}

\maketitle

\begin{abstract}
视觉的根本目的并不是输出一串标签，而是从不完整、受噪声污染的输入中恢复一个足够稳定、可操作、可用于定位、预测与想象的外部世界表征。受心理物理学、系统神经科学与 psychometrics 中“通过密集行为探测反推内部表征结构”的思路启发，本文提出 \textbf{Ordinary-Bench}：一个通过受控 3D 场景、序数空间关系问答与场景级重建来评估视觉语言模型（VLM）空间信念的 benchmark。本文的核心观察是：既然充分且准确的偏序关系可以近似重建场景，那么对 VLM 进行大规模 QRR/TRR 探测，本质上就是在读取它“脑海中的场景”是否存在、是否一致、是否可复原。基于当前 540 个场景、200000+ QA 对的实验，我们发现模型在空间信念层面的能力远弱于表面问答分数所暗示的水平：单图条件下 QRR 准确率约为 70\%，多视角输入反而下降到接近随机猜测，TRR 长期停留在随机水平，而基于 200K+ QA 对的微调或强化学习也没有带来显著改善。与生成式 world model 或依赖工具使用、长思维链和 test-time search 的 visual CoT 方法不同，Ordinary-Bench 评估的不是模型能否借助额外流程把题做对，而是其回答本身能否定义一个对象级、代码级、可矢量化重建的场景结构。由此，我们将评测问题从“模型答对了多少空间题”推进到一个更严格的问题：\emph{模型是否真正形成了一个可被重建、可被想象、可被对齐到真实世界的空间场景信念}，并为后续的场景矢量化重建、定位与结构化视觉记忆提供可操作的评估路径。

\end{abstract}

\section{引言}

人类视觉系统的价值，并不只体现在“认出了什么”，更体现在“构成了怎样的世界”。从视网膜上的二维投影到对物体位置、朝向、距离、遮挡、可达性与未来状态的判断，视觉的深层功能一直是把局部、瞬时、噪声化的感受，转化为一个足够稳定的外部场景模型。正因为如此，在心理物理学、系统神经科学与 psychometrics 中，研究者很少满足于单个刺激上的正确率；他们更关心观察者在系统性探测、重复测量、条件扰动和跨任务比较下，是否表现出一致的结构、稳定的表征与可解释的失真模式 \cite{kriegeskorte2008rsa,yamins2016dicarlo,xu2025bsa,zhang2025comfort}。如果我们真的希望理解 VLM “看见了什么”，那么就不该只把它当作一个会答题的黑箱，而应该像研究一个观察者那样，去追问：它的回答背后是否对应着一个足够一致、足够可重建的世界。

现有 VLM 评测与这一目标之间存在明显落差。大量工作把空间理解压缩成单题问答、选择题正确率或局部关系判断：模型说对了“左边”“更近”“前方”这样的答案，就被视为具备某种空间理解能力。然而，答对局部问题，并不等于形成了全局场景；多个单独正确的回答，也完全可能彼此不兼容。一个模型也许能在单图条件下猜对一部分距离比较，却无法把这些判断拼接成一致的布局；它也可能在语言模板的帮助下给出看似合理的方位答案，但这些答案一旦进入统一坐标系，就会立即暴露出冲突、塌缩和不可实现性。换句话说，当前许多 benchmark 测到的，其实更像是“局部关系回答能力”，而不是“可被重建的空间信念”。

本文的核心动机来自序数嵌入与场景重建的一个基本事实：当偏序关系足够充分且足够准确时，场景几何可以在相似变换意义下被近似恢复 \cite{agarwal2007gnmds,ma2018robustordinal,dokmanic2015edm}。这意味着，如果我们不再满足于向 VLM 提几个零散问题，而是系统性地询问大量相对距离与相对方位关系，那么模型的回答集合本身就构成了一种行为学意义上的“断层扫描”。在这个视角下，QRR 与 TRR 不只是两个问答任务，而是两类读取模型外显场景信念的探针：QRR 通过距离偏序检验模型是否保留了物体间的相对近远结构；TRR 通过方向关系检验模型是否维持了统一的参照系与方位几何。若这些回答能够支撑一个与真实场景接近的重建结果，我们才有理由说模型确实形成了某种场景级表征；反之，如果这些回答只能生成一个高度冲突、严重失真或完全不可实现的布局，那么模型即使在若干题目上答对了，也更可能只是依赖语言先验、模板匹配或彼此独立的局部启发式。

这也使 Ordinary-Bench 与两类看似相近、但目标不同的研究形成了清楚区分。第一类是 world model 工作：从经典的 latent dynamics 建模，到近期用于导航、驾驶和 4D 场景预测的生成式模型，它们的核心目标通常是预测未来观测、生成视频、BEV 或 occupancy，并服务于规划与控制 \cite{ha2018worldmodels,wang2024drivewm,bar2025navigationwm,ni2025maskgwm,liao2025i2world,zhou2025hermes}。第二类是 visual CoT / tool-use 路线：模型通过外部视觉工具、代码执行、视觉草图板、树搜索或更长的 test-time reasoning 来把题做对 \cite{suris2023vipergpt,liu2023llavaplus,gui2024vpd,hu2024visualsketchpad,wang2025visuothink,lin2025multimodalcot}。这两条路线都很重要，但它们优化的主要是生成、规划或解题流程，而不是回答本身是否已经定义了一个对象级、可反演、可矢量化的场景结构。Ordinary-Bench 要测的，正是这一更基础的问题。

基于这一思路，我们提出 \textbf{Ordinary-Bench}。它不是再增加一个“空间题库”，而是把 VLM 当作一个需要被行为反演的观察者：我们在受控 3D 渲染环境中生成结构明确的场景，以 QRR（四元相对关系）和 TRR（三元时钟方向关系）进行密集探测，再将模型输出转写为偏序约束与方向约束，检验这些回答是否足以支持一个一致的场景重建。使用受控合成环境而不直接依赖自然图像，是一个刻意的选择。对于“模型究竟构成了什么场景”这一问题，精确几何真值、可控复杂度和可重复扰动，比生态真实性更重要；只有在场景结构本身无歧义的条件下，我们才能严格地区分视觉贡献、语言先验和结构性失真 \cite{johnson2017clevr,yi2020clevrer}。

当前已经完成的实验结果进一步说明了为什么需要这种更严格的评估框架。基于目前整理出的 540 个场景与 200000+ QA 对，我们观察到：在单图条件下，模型在 QRR 上还能达到大约 70\% 的准确率，看起来似乎具备了某种距离比较能力；但一旦输入切换到多视角，QRR 表现反而退化到接近随机猜测，表明模型并没有把多视角信息整合为更稳定的空间表征，而更像是在额外视角的干扰下丧失了原本脆弱的局部线索。TRR 的情况更差，几乎始终停留在随机猜测水平。这意味着当前 VLM 对“谁更近”可能还有有限的、脆弱的判断能力，但对“统一参照系中的相对方向”几乎没有形成可靠表征。更值得注意的是，基于 540 个场景、200000+ QA 对进行微调或强化学习，并没有带来显著的提高和优化。这一负结果具有重要含义：仅仅把更多关系问答灌输给模型，并不会自动产生几何一致的场景信念；如果评测目标仍停留在局部正确率上，这种结构性失败很容易被掩盖。

从更大的应用角度看，这一问题并不只是“空间题有没有答对”。如果一个 VLM 真的能在偏序关系层面形成可重建、可想象、可对齐的场景信念，那么它实际上已经具备了一种与具体数值尺度弱耦合的场景矢量化能力。偏序重建并不要求模型先恢复每一条精确欧氏距离；只要关系结构稳定存在，我们就可以在后处理中再引入真值尺度、经验尺度或任务相关尺度，把相对布局对齐为可用于定位、检索、导航、交互甚至场景编辑的结构化表示。也就是说，Ordinary-Bench 评估的并不是一个狭义的 QA 技能，而是在测试 VLM 是否已经触碰到“从视觉中恢复可操作场景结构”这一更本质的能力。

进一步地，Ordinary-Bench 不应只停留在静态场景。既然真实视觉的目标是维持一个随时间更新的场景世界，那么更合理的实验设计应该同时覆盖：单图下的局部关系感知，多视角下的统一场景整合，以及动态场景中关系随时间变化后的持续一致性。也就是说，未来实验不只应问“当前谁更近、谁在谁的左边”，还应问“何时发生关系翻转、何时出现交叉或遮挡、模型在时间推进中是否仍维持同一个场景解释”。这一点也与近期视频与 world-model 评测中对 motion-level、temporal consistency 和 future-state prediction 的关注形成呼应 \cite{hong2025motionbench,fu2024videomme,zhou2024mlvu,liu2024tempcompass,gao2025wmabench}。

本文的贡献可以概括为四点。第一，我们提出一个受心理物理学启发的 VLM 空间评估框架，用受控 3D 场景和密集 QRR/TRR 探测，把模型从“答题器”重新定义为“可被行为反演的观察者”。第二，我们将序数嵌入中的可重建性思想引入 VLM 评测：通过足够多的偏序问答，读取并检验模型的外显场景信念，而不再只报告单题正确率。第三，我们给出一组当前已经相当清晰的负面证据：单图下有限的 QRR 成功并不意味着模型形成了稳定空间表征；多视角并未提升而是压垮了这种能力，TRR 长期接近随机，且微调与强化学习并未实质改变这一局面。第四，我们进一步提出一个从静态到动态的实验设计方向，把多视角整合、时序关系变化、场景级可实现性与工具增强对照纳入统一 protocol。综合来看，本文试图把问题从“VLM 会不会回答空间问题”推进到一个更根本也更严格的问题：\emph{VLM 是否真的在视觉输入上构造了一个足够一致、足够可复原的世界？}

\section{相关工作}

\subsection{心理物理学、系统神经科学、psychometrics 与行为反演}

本文的方法论首先受到心理物理学、系统神经科学与 psychometrics 的启发。在这些传统中，研究对象往往不是通过一次性的“答对/答错”被定义，而是通过大量受控刺激、重复测量、条件扰动和跨任务比较来刻画其内部表征的结构。表征相似性分析（RSA）强调，可以不直接观测底层机制，而通过系统的行为或响应模式来反推表征空间的组织方式 \cite{kriegeskorte2008rsa}；目标驱动模型理解感觉皮层的工作进一步表明，真正重要的并不只是单任务得分，而是模型是否复现了系统性结构与稳定不变性 \cite{yamins2016dicarlo}。近期面向 VLM 的 psychometric 工作开始进一步拆解空间能力本身：Basic Spatial Abilities 将空间感知、关系、定向、心像旋转和空间可视化拆成多个基本能力维度 \cite{xu2025bsa}，COMFORT 则系统研究了空间语言中 frame of reference 的歧义、一致性与跨语言差异 \cite{zhang2025comfort}。本文借用的是这种“行为可检验、结构可反演、能力可分解”的研究姿态：我们并不声称直接读出 VLM 的内部神经状态，而是把大量 QRR/TRR 回答视为其外显空间信念的可观测痕迹。

\subsection{通用多模态 benchmark 与评测错配}

近两年的多模态 benchmark 一方面大幅拓宽了任务覆盖，另一方面也暴露出“高分不等于真实视觉理解”的问题。MMMU 与 SEED-Bench 从大规模、多学科、多图像类型的角度建立了广谱 multimodal evaluation 版图 \cite{yue2024mmmu,li2024seedbench}，但随后的 MMMU-Pro 进一步证明：许多题目在 text-only 设定下依然可解，且 benchmark shortcut 可能显著高估模型的视觉能力 \cite{yue2025mmmupro}。另一条线则更直接地质疑现有评测是否测到了它们声称要测的东西。HallusionBench 通过控制组设计剖析语言 hallucination 与 visual illusion \cite{guan2024hallusionbench}；BLINK 强调模型“能看见但不一定真正感知” \cite{fu2024blink}；MMStar 把视觉非必要性和评测泄漏显式写成 benchmark 设计问题 \cite{chen2024mmstar}；POPE 和 VLInd-Bench 则分别从目标幻觉与语言先验污染的角度暴露了当前评测范式的脆弱性 \cite{li2023pope,lee2025vlind}。Ordinary-Bench 与这些工作同向之处在于都不满足于表面正确率；不同之处在于，我们进一步把评价终点推进到“回答集合是否共同定义了一个真实、可实现、可重建的场景”。

\subsection{空间、三维、多视角、遮挡与动态 benchmark}

空间理解 benchmark 本身也在快速细化。VSR、What'sUp 与 SpatialVLM 指出，哪怕是看似基础的左右、上下、前后与空间关系，VLM 也远没有真正掌握 \cite{liu2023vsr,kamath2023whatsup,chen2024spatialvlm}。随后的一批工作开始追求更干净或更细分的空间评测：SpatialMQA 尝试避免 bbox shortcut 与纯先验作答 \cite{liu2025spatialmqa}；SPHERE 以层级化方式从单技能、多技能整合到高层空间推理系统测量空间盲点 \cite{zhang2025sphere}；Can Transformers Capture Spatial Relations Between Objects? 强调应采用物理上更明确的关系定义，而不是含混的语言学关系类别 \cite{wen2024spatialrelations}。再往前一步，3DSRBench 将问题推进到自然图像中的 3D awareness、视角变化和 FlipEval 稳健性 \cite{ma2025threedsrbench}；CAPTURE 从遮挡与 amodal counting 角度测试模型能否在缺失视觉证据时维持结构想象 \cite{pothiraj2025capture}；MetaVQA 与 NuPlanQA 则把多视角、ego-centric、embodied driving scenes 纳入 benchmark \cite{wang2025metavqa,park2025nuplanqa}。在时间维度上，Video-MME、MLVU、TempCompass 与 MotionBench 进一步表明，视频 VLM 在 temporal reasoning 与最基础的 motion-level perception 上同样存在显著缺口 \cite{fu2024videomme,zhou2024mlvu,liu2024tempcompass,hong2025motionbench}。这些工作都在逼近 Ordinary-Bench 所关心的问题，但它们大多仍以题目级、任务级或 skill-level 得分为中心；本文进一步追问的是：这些回答能否被拼接成一个统一场景。

\subsection{对象级场景表示、amodal completion 与 world model}

另一条与本文相邻的路线关注对象级场景结构与世界模型。CObL 利用 object layers 表示图像中的遮挡顺序与 amodal 物体层，强调对象级、分层化的场景组织 \cite{damaraju2025cobl}；CAPTURE 也可被视为测试模型是否具有在遮挡下补全隐藏结构的能力 \cite{pothiraj2025capture}。更广义地，从经典的 World Models 到近期的 Drive-WM、Navigation World Models、MaskGWM、I2-World 与 HERMES，world model 文献的核心目标通常是学习潜在动力学并预测未来像素、视频、BEV 或 occupancy 演化，以服务于规划、控制和仿真 \cite{ha2018worldmodels,wang2024drivewm,bar2025navigationwm,ni2025maskgwm,liao2025i2world,zhou2025hermes}。与此同时，也开始出现“模型到底有没有 coherent world model”这一更严格的评测工作：WM-ABench 把世界模型能力拆成 perception 与 prediction 两阶段进行原子化评估 \cite{gao2025wmabench}；Vafa 等则进一步指出，表面上的单步预测成功并不意味着模型真的学到了连贯的内部状态结构 \cite{vafa2024implicitwm}。本文与这两条线都相关，但又有明显区别：我们既不训练一个生成式 world model，也不要求模型输出像素、视频或 occupancy；我们评估的是回答本身是否定义了一个对象级、关系级、代码级、可矢量化的场景结构。

\subsection{视觉 CoT、工具使用与 test-time scaling}

近期另一条非常活跃的方向是 visual CoT、工具使用与 inference-time scaling。ViperGPT 通过代码生成显式组合检测、深度、计数等视觉工具 \cite{suris2023vipergpt}；LLaVA-Plus 将多模态 agent 与技能库连接起来 \cite{liu2023llavaplus}；Visual Program Distillation 则进一步把“工具 + 程序化推理”的能力蒸馏进单一 VLM \cite{gui2024vpd}。与此同时，Visual Sketchpad 让模型通过画线、框和标记把中间视觉推理外化 \cite{hu2024visualsketchpad}；VisuoThink 和 multimodal CoT inference-time scaling 则把树搜索、更长的 test-time reasoning 与视觉中间状态结合起来 \cite{wang2025visuothink,lin2025multimodalcot}。这些方法能够显著提升题目正确率，但它们优化的主要是“如何借助额外流程把题做对”，而不是“模型原本是否已经形成稳定的视觉场景表征”。因此，在本文中，这些工作更适合作为能力增强路线与目标错位对照：即便它们能提升分数，也不自动说明模型本身已经拥有一个可重建的场景世界。

\subsection{可靠性、选择性回答与几何可实现性}

当回答会被解释为场景约束时，可靠性问题会变得更加关键。Reliable VQA 提出，在证据不足时拒答往往比错误作答更可靠 \cite{whitehead2022reliablevqa}；选择性回答与一致性评测工作进一步把不确定性纳入黑箱 VLM 评价 \cite{eisenschlos2024selectivevqa,khan2024consistency}。而在理论层面，本文最直接的基础来自序数嵌入与几何可实现性：Generalized Non-metric MDS 和后续鲁棒序数嵌入工作说明，基于相对距离大小而非绝对距离值，同样可以恢复点集结构 \cite{agarwal2007gnmds,ma2018robustordinal}；Euclidean Distance Matrix 理论提醒我们，约束集的存在并不自动意味着其对应某个可实现的欧氏场景 \cite{dokmanic2015edm}；对于 TRR 一类方向关系，bearing rigidity 和 angle rigidity 则提供了关于唯一性、镜像歧义与方向约束可实现性的几何语言 \cite{zhao2019bearingrigidity,chen2021anglerigidity}。Ordinary-Bench 的创新不在于提出新的嵌入算法，而在于把这些理论转化为一种 VLM 评测终点：既然充分准确的偏序关系可以支撑场景重建，那么对 VLM 进行密集关系问答，本质上就是在测试它是否拥有一个足以支持重建的空间信念。

\subsection{与本文的区别}

综合来看，现有工作大致分布在五个方向：一类研究 VLM 是否会做广义多模态题，一类研究空间与 3D / 多视角 / 动态题能否答对，一类研究评测是否被语言先验和 benchmark shortcut 污染，一类研究 world model 的生成与预测能力，另一类研究是否能通过工具与更长推理链提升最终得分。本文试图把这些线索收束到一个更严格的判断标准上。我们既继承心理物理学式的黑箱观察者视角，也采用受控合成世界的实验设计，同时用序数重建的可实现性作为评测终点，并明确把“对象级、代码级、可矢量化的场景结构”与生成式 world model 或工具增强解题流程区分开来。由此，Ordinary-Bench 不是单纯增加一个空间 benchmark，而是提出一种新的判断标准：一个 VLM 只有在其大量关系回答能够共同支持一个一致、可复原、可对齐的场景时，才算真正“看见了”空间。


\section{Ordinary-Bench：数据与协议}
该节后续将继续在本目录中扩写，统一迁移到官方模板工作流下。

\section{场景信念重建}
该节后续将继续在本目录中扩写，重点围绕 QRR/TRR 约束到场景外显信念的行为反演。

\section{实验设计与当前观察}
\subsection{研究问题}

本文的实验设计围绕五个问题展开。第一，VLM 在单图条件下是否真的具备稳定的局部空间关系理解能力，而不是只会利用语言模板或局部纹理线索。第二，多视角输入是否会把单视角中的脆弱正确率整合为更稳定的场景表征，还是会暴露模型缺乏统一场景信念的问题。第三，模型给出的 QRR/TRR 回答能否被联合解释为一个可实现、可重建、与真实场景对齐的对象级结构。第四，这种结构一旦放入时间维度，在物体运动、遮挡和关系翻转出现后能否继续保持一致。第五，当引入 CoT、visual CoT、工具使用或 test-time search 后，模型提升的究竟是“题目得分”，还是“场景级可重建性”。

\subsection{静态场景协议}

静态实验建议至少包含三层设置。第一层是 \textbf{单图局部关系评测}：对每个场景采样大量 QRR 和 TRR 问题，分别统计 item accuracy，并按关系类型、对象类别数、遮挡强度和视角复杂度分桶。第二层是 \textbf{多视角整合评测}：向模型提供同一场景的多个观察视角，要求其回答同样的 QRR / TRR 集合，并显式比较单图到多视角的性能变化。第三层是 \textbf{场景级重建评测}：把模型回答转写为距离偏序与方向约束，在统一坐标系中求解一个近似重建布局，再测量其可实现率、冲突率、重建误差、镜像歧义情况以及与真实场景的一致性。

在指标上，单题准确率只能作为最低层统计，真正关键的是 scene-level 指标。建议至少报告：
\begin{itemize}[leftmargin=1.5em]
\item \textbf{Constraint realizability}: 回答集合是否存在统一几何解释。
\item \textbf{Contradiction rate}: 约束之间是否出现环、互斥或跨视角冲突。
\item \textbf{Reconstruction error}: 重建布局与真实场景在相似变换对齐后的结构误差。
\item \textbf{View-consistency gap}: 单图与多视角条件下 scene-level 指标的差值。
\item \textbf{Belief collapse rate}: 多视角输入后从“局部可答”退化为“整体不可实现”的比例。
\end{itemize}

\subsection{当前静态观察}

基于当前已经完成的 540 个场景和 200000+ QA 对，我们已经得到一组对论文中心论点非常关键的初步观察。第一，在单图条件下，QRR 准确率约为 70\%，说明模型对局部相对远近并非完全无感，但这种能力仍明显弱于一个真正稳定的空间系统。第二，当输入扩展为多视角时，QRR 表现并未提升，反而下降到接近随机猜测，这表明模型没有把多个视角整合为更完整的场景表征，而更像是被额外视觉证据破坏了原本脆弱的局部启发式。第三，TRR 长期停留在随机水平，说明模型对统一参照系中的相对方向关系几乎没有形成可靠表征。第四，基于 200K+ QA 对进行微调或强化学习并没有带来显著改善，意味着单纯增加关系监督并不会自动产生几何一致的场景信念。

\subsection{动态场景扩展}

如果论文要从“视觉的根本目的”这一更大叙事切入，那么实验设计不应停留在静态图片。更合理的下一步是引入 \textbf{动态动画场景}，让 benchmark 从静态空间扩展到时空一致性。当前仓库中已有动态图生成与时序关系抽取管线，这为后续实验提供了很好的实现基础，但实验设计本身不需要被现有代码完全限制。

动态场景可以围绕四类任务设计。第一类是 \textbf{时序 QRR}：在不同时间步上询问对象对之间的相对远近，观察模型是否能维持随轨迹变化而更新的偏序结构。第二类是 \textbf{时序 TRR}：询问相对方向关系在运动过程中的演化，测试模型能否维持统一参照系并追踪角度变化。第三类是 \textbf{change-point probing}：直接提问关系何时发生翻转、接近、交叉、相遇、遮挡或离开。第四类是 \textbf{trajectory-consistency probing}：分别根据单帧、短序列、长序列和多视角序列提问，比较模型隐含的场景重建是否在时间推进中仍然相容。

动态实验的核心不应只是“视频题能不能答对”，而应测量模型是否维持了一个 \textbf{随时间更新但结构连续的世界}。因此建议新增以下指标：
\begin{itemize}[leftmargin=1.5em]
\item \textbf{Temporal consistency}: 相邻时间步重建结果的平滑性与因果连续性。
\item \textbf{Relation change accuracy}: 对关系翻转、接近、远离、遮挡事件的识别准确率。
\item \textbf{Future belief stability}: 让模型根据前若干帧预测未来关系，再与真值轨迹对比。
\item \textbf{Cross-time realizability}: 把多个时间步的回答联立后，是否仍存在统一的时空解释。
\end{itemize}

\subsection{干预、对照与能力归因}

为了区分“模型本身具备场景信念”与“模型借助额外流程把题做对”，实验中建议设置两大类 protocol。第一类是 \textbf{原生能力 protocol}：直接回答、简短 CoT、标准多视角输入。第二类是 \textbf{增强能力 protocol}：加入 visual CoT、草图板、程序执行、外部检测器 / 分割器、树搜索或更长的 inference-time scaling \cite{hu2024visualsketchpad,wang2025visuothink,lin2025multimodalcot,suris2023vipergpt,gui2024vpd}。如果增强方法只提升 item accuracy 而不能显著改善 scene-level reconstructability，那么就应将其解释为“解题流程增强”，而不是“真实视觉表征增强”。

\subsection{论文中建议出现的核心表图}

实验部分建议至少包含四类主表。第一类是单图与多视角条件下的 QRR / TRR 准确率表，直接展示从单图到多视角的退化。第二类是 scene-level 指标表，包括 realizability、contradiction rate 与 reconstruction error。第三类是训练干预表，对比 zero-shot、few-shot、SFT、RL 以及 tool-use / visual CoT 方案。第四类是动态实验表，展示 temporal consistency、relation change accuracy 与 cross-time realizability。与之配套的图示建议包括：多视角导致场景信念塌缩的可视化例子、重建成功与失败的场景对比图、以及动态场景中关系随时间翻转的轨迹图。


\section{讨论与应用展望}
该节后续将继续在本目录中扩写，重点讨论局部答对与全局可重建之间的差异，以及偏序驱动的场景矢量化表示前景。

\bibliographystyle{unsrtnat}
\bibliography{spatiotemporal_vlm_refs}

\end{document}
